% page 325 of the facsimile (image p-324 of the scan)
% block 157.5 x 246.3 mm ; left margin 44.5 mm
\pgc{22.860mm}{0.000mm}{
\l{\cel{52}317}
\l{}
\l{kromatino. (\sou{biol.})\cel{2}Substanco tre afina pri la kolorizivi}
\l{\cel{3}bazala, e quan onu renkontras maxim-multa-kaze forme}
\l{\cel{3}di grani kunjuntita a la filamento nukleala.}
\l{}
\l{kromatofila. (\sou{biol.)}\cel{4}Dicesas pri la celuli, grani od}
\l{\cel{3}+organiti, apta kaptar facile la kolorizivi.}
\l{}
\l{kromatoforo.\cel{2}(\sou{biol.}) Celulo di la tisuo konjuntiva, qua}
\l{\cel{3}kontenas multa pigmenti, quan onu renkontras periferie di}
\l{\cel{3}la neuraxo e di la membrani sensala (iriso, pelo, ec.)}
\l{}
\l{\sou{kromatolizo.}\cel{2}(\sou{biol.)}\cel{2}Dissolveso di la kromatino di la celu-}
\l{\cel{3}lo-nukleo,qua karakterizas un ek la fazi di mitoso.}
\l{}
\l{\sou{kromidio}. (\sou{biol.})\cel{2}Nomo donita da Hertwig a la elementi di}
\l{\cel{3}la nukleo nepreciza di la bakterii.}
\l{}
\l{\sou{kromolitografar}. (\sou{trans.}) Agar impro litografiala koloroza.}
\l{}
\l{\sou{kromosomo}. (\sou{bio}l.)\cel{2}Segmento di la filamento kromatika nu-}
\l{\cel{3}klea qua aparas dum mitoso.}
\l{}
\l{\sou{kromotipar}. (\sou{trans.}) Imprimar koloroze per la procedi tipo-}
\l{\cel{3}grafiala. - DEFRS.}
\l{}
\l{\sou{\textquotedbl{}kron\textquotedbl{}}. Monet-uno, di qua la valoro varias segun la landi}
\l{\cel{3}(Britania, Suedia, e c.)}
\l{}
\l{\sou{krono}. I.Ringo destinita a metesar cirkum la kapo : (1) Rin-}
\l{\cel{3}go de flori, de folii interplektita quan onu metas cirkum}
\l{\cel{3}la kapo kom (a) ornivo ; (b) distingivo armeala o civila.}
\l{\cel{3}(\sou{metaf.}) Rekompenso cielala) ; (2) Ringo de metalo, orni-}
\l{\cel{3}ta ye floroni, ye perli, e c. quan onu metas cirkum la}
\l{\cel{3}kapo kom insigno di la suvereneso, di la nobeleso. La}
\l{\cel{3}krono rejala, dukala, e c. qua esas super la stampuro}
\l{\cel{3}blazonala. - II. To di quo la formo memorigas krono. -}
\l{\cel{3}III.\cel{1}\sou{Dento-krono}\cel{1}: la parto supra qua salias de la jen-}
\l{\cel{3}jivo e kovresas ye emalio. - DEFIRS.}
\l{}
\l{\sou{kronglaso}. Vitro diafana di tre bela qualeso,qua kontenas}
\l{\cel{3}multa potaso, tre facile gisebla (fuzebla), di qua la}
\l{\cel{3}refraktiveso-grado esas tre alta, e qua uzesas por fabri-}
\l{\cel{3}kar lensi por la instrumenti optikala, ed anke prismi kro-}
\l{\cel{3}mata.DEFIRS.}
\l{}
\l{\sou{kronika}. (\sou{aludante morbo)}\cel{2}Qua par-iras lente sua periodi.}
\l{\cel{3}II. Qua divenis forta segun ke tempo fluis e, konseque,}
\l{\cel{3}ne plus esas akuta. - DEFIRS.}
\l{}
\l{\sou{kroniko}. I. Kolektajo de fakti historiala, segun la ordino}
\l{\cel{3}di\cel{1}\sou{lia}\cel{1}evento. - II. Ensemblo ek la fakti, ek la informi}
\l{\cel{3}qui esas mencionita pri la personi. - DEFIRS.}
\l{}
\l{kronologio. La cienco di qua la temo esas pruvar la dati di la}
\l{\cel{3}eventaji.historiala. - DEFIRS.}
}
